\subsection{论文结构及各部分写法}

\small
\begin{longtable}{p{2.25cm}p{4.0cm}p{5.1cm}p{4.2cm}}
\caption{A053章节功能与本队理解}\label{tab:a053-structure}\\
\toprule
部分 & 论文在写什么 & 具体写法 & 本队吸收的做法\\
\midrule
\endfirsthead
\toprule
部分 & 论文在写什么 & 具体写法 & 本队吸收的做法\\
\midrule
\endhead
\bottomrule
\endfoot
摘要（p.1）
& 依次回答运动状态、碰撞、最小螺距、调头曲线和速度上限。
& 每个分问都给对象、模型动作和交付值，如第一次碰撞时刻、螺距和龙头速度；算法名放在它所求的量之后。
& 摘要按分问写，每问至少留下一个有判别力的结果；减少重复的顺序词，不把摘要写成章节目录。\\

问题分析（p.3--4）
& 把板凳龙的整体运动拆成相邻把手位置递推，再说明各问增加的约束。
& 问题四被拆成“求较短调头曲线”和“求调头过程中的运动状态”两个子任务，前者输出路径，后者沿用位置递推。
& 先画清输入输出关系，再选方法；复用前问时点明沿用量和新增量。\\

模型假设（p.2）
& 限定人为扰动、板凳刚体和开始调头的时机。
& 假设会进入方程或事件判定，没有写与计算无关的常识句。
& 每条假设说明“它使哪一项可以怎样计算”，并在评价中说明放宽后要改哪个模块。\\

符号说明（p.3）
& 给出螺距、极角、位置、速度、圆弧半径等贯穿五问的量。
& 统一公共符号，分问新增量在首次使用处定义，减少同一字母承担两个含义。
& 符号表只放多次使用的量；仅在一处使用的中间量随公式解释。\\

模型建立（p.6--21、p.23--36）
& 写螺线弧长、定长递推、矩形碰撞、相切圆弧和多段路径状态。
& 公式按“已知把手--定长方程--候选根--运动方向筛选--下一把手”推进；碰撞从几何对象落到矩形投影区间。
& 让公式顺序接近程序顺序。遇到多根时写清筛根条件，遇到分段路径时先列状态和切换条件。\\

模型求解（p.8--16、p.22--31、p.34--41）
& 用二分、时间离散、变步长遍历、状态递推和粒子群得到五问结果。
& 螺距先在 $[45,46]\,\mathrm{cm}$ 内定位，再以 $0.01\,\mathrm{cm}$ 步长细查；碰撞检测在第一次触发时停止。
& 粗查回答“边界在哪一段”，细查回答“边界精确到多少”；随机算法另报种子和重复次数。\\

结果及分析（p.22、p.27--28、p.32--33、p.38、p.41）
& 报告碰撞时刻、螺距、圆弧半径、曲线长度和速度上限，并解释局部异常。
& 参数微变后若第一次碰撞对象改变，论文从局部几何解释输出跳变；12--13秒的速度突增则用相同时间内的路程差解释。
& 表后不复述数字，改为回答“为何变、哪条约束起作用、结论只到什么范围”。\\

检验与边界
& 检查首次碰撞对象、全过程可行性和现实动作限制。
& p.17用反设和逐节前推说明最早碰撞对象；p.28用“终点可行但中途碰撞”的方案否定只查终点；p.32从龙身不能倒退推出半径下界。
& 检验要对应真实疑问。回代、反例和物理边界各自有用途，不能合写成一句“模型得到验证”。\\

附录（p.43--62）
& 给出逐把手循环、二分、分段路径、碰撞和粒子群代码。
& 正文四种运动状态与代码状态码1--4相互对应，表格输出也由固定变量名生成。
& 附录保留能复现正文结果的入口、参数和输出说明；正文只留关键流程，不粘贴整段程序。\\
\end{longtable}
\normalsize

